%%%%%%%%%%%%%%%%%%%%%%%%%%%%%%%%%%%%%%%%%
% Wenneker Assignment
% LaTeX Template
% Version 2.0 (12/1/2019)
%
% This template originates from:
% http://www.LaTeXTemplates.com
%
% Authors:
% Vel (vel@LaTeXTemplates.com)
% Frits Wenneker
%
% License:
% CC BY-NC-SA 3.0 (http://creativecommons.org/licenses/by-nc-sa/3.0/)
% 
%%%%%%%%%%%%%%%%%%%%%%%%%%%%%%%%%%%%%%%%%

%----------------------------------------------------------------------------------------
%	PACKAGES AND OTHER DOCUMENT CONFIGURATIONS
%----------------------------------------------------------------------------------------

\documentclass[11pt]{scrartcl} % Font size

%%%%%%%%%%%%%%%%%%%%%%%%%%%%%%%%%%%%%%%%%
% Wenneker Assignment
% Structure Specification File
% Version 2.0 (12/1/2019)
%
% This template originates from:
% http://www.LaTeXTemplates.com
%
% Authors:
% Vel (vel@LaTeXTemplates.com)
% Frits Wenneker
%
% License:
% CC BY-NC-SA 3.0 (http://creativecommons.org/licenses/by-nc-sa/3.0/)
% 
%%%%%%%%%%%%%%%%%%%%%%%%%%%%%%%%%%%%%%%%%

%----------------------------------------------------------------------------------------
%	PACKAGES AND OTHER DOCUMENT CONFIGURATIONS
%----------------------------------------------------------------------------------------

\usepackage{amsmath, amsfonts, amsthm} % Math packages
\usepackage{listings} % Code listings, with syntax highlighting
\usepackage[english]{babel} % English language hyphenation

\usepackage{graphicx} % Required for inserting images
\graphicspath{{Figures/}{./}} % Specifies where to look for included images (trailing slash required)

\usepackage{booktabs} % Required for better horizontal rules in tables

% \numberwithin{equation}{section} % Number equations within sections (i.e. 1.1, 1.2, 2.1, 2.2 instead of 1, 2, 3, 4)
\numberwithin{figure}{section} % Number figures within sections (i.e. 1.1, 1.2, 2.1, 2.2 instead of 1, 2, 3, 4)
\numberwithin{table}{section} % Number tables within sections (i.e. 1.1, 1.2, 2.1, 2.2 instead of 1, 2, 3, 4)

\setlength\parindent{0pt} % Removes all indentation from paragraphs

\usepackage{enumitem} % Required for list customisation
\setlist{noitemsep} % No spacing between list items

%----------------------------------------------------------------------------------------
%	DOCUMENT MARGINS
%----------------------------------------------------------------------------------------

\usepackage{geometry} % Required for adjusting page dimensions and margins

\geometry{
	paper=a4paper, % Paper size, change to letterpaper for US letter size
	top=2.5cm, % Top margin
	bottom=3cm, % Bottom margin
	left=3cm, % Left margin
	right=3cm, % Right margin
	headheight=0.75cm, % Header height
	footskip=1.5cm, % Space from the bottom margin to the baseline of the footer
	headsep=0.75cm, % Space from the top margin to the baseline of the header
	%showframe, % Uncomment to show how the type block is set on the page
}

%----------------------------------------------------------------------------------------
%	FONTS
%----------------------------------------------------------------------------------------

\usepackage[utf8]{inputenc} % Required for inputting international characters
\usepackage[T1]{fontenc} % Use 8-bit encoding

\usepackage{fourier} % Use the Adobe Utopia font for the document

%----------------------------------------------------------------------------------------
%	SECTION TITLES
%----------------------------------------------------------------------------------------

\usepackage{sectsty} % Allows customising section commands
\sectionfont{\vspace{6pt}\normalfont\scshape} % \section{} styling
\subsectionfont{\normalfont\bfseries} % \subsection{} styling
\subsubsectionfont{\normalfont\itshape} % \subsubsection{} styling
\paragraphfont{\normalfont\scshape} % \paragraph{} styling

%----------------------------------------------------------------------------------------
%	HEADERS AND FOOTERS
%----------------------------------------------------------------------------------------

\usepackage{scrlayer-scrpage} % Required for customising headers and footers

\ohead*{} % Right header
\ihead*{} % Left header
\chead*{} % Centre header

\ofoot*{} % Right footer
\ifoot*{} % Left footer
\cfoot*{\pagemark} % Centre footer
 % Include the file specifying the document structure and custom commands

%----------------------------------------------------------------------------------------
%	TITLE SECTION
%----------------------------------------------------------------------------------------
% Add a logo to the title page from the folder "Figures"





\title{	
	\includegraphics[width=0.25\textwidth]{./Figures/Sharif_University_Logo.jpg}\\
	\normalfont\normalsize
	\textsc{Sharif University of Technology}\\ % Your university, school and/or department name(s)
	\vspace{25pt} % Whitespace
	\rule{\linewidth}{0.5pt}\\ % Thin top horizontal rule
	\vspace{20pt} % Whitespace
	{\huge Deep Learning Assignment 3}\\ % The assignment title
	\vspace{12pt} % Whitespace
	\rule{\linewidth}{2pt}\\ % Thick bottom horizontal rule
	\vspace{12pt} % Whitespace
}
% \includegraphics[width=0.1\textwidth]{C:/Users/saeedzou/Documents/sharif_logo.png}~[1cm]

\author{\LARGE Saeedreza Zouashkiani} % Your name
\DeclareMathOperator{\E}{\mathbb{E}}
\DeclareMathOperator*{\argmin}{argmin}

\date{\normalsize\today} % Today's date (\today) or a custom date

\begin{document}
% \titlegraphic{\includegraphics[width=0.5\textwidth]{./Figures/sharif_logo.png}}
\maketitle % Print the title

\section{} %1
\subsection{} %1.1
The output of the following image
\begin{equation}
\begin{bmatrix}
1 & 0 & 1 & 0 & 1 & 0 & 1 & 0\\
0 & 1 & 0 & 1 & 0 & 1 & 0 & 1\\
1 & 0 & 1 & 0 & 1 & 0 & 1 & 0\\
0 & 1 & 0 & 1 & 0 & 1 & 0 & 1\\
1 & 0 & 1 & 0 & 1 & 0 & 1 & 0\\
0 & 1 & 0 & 1 & 0 & 1 & 0 & 1\\
1 & 0 & 1 & 0 & 1 & 0 & 1 & 0\\
0 & 1 & 0 & 1 & 0 & 1 & 0 & 1\\
\end{bmatrix}
\end{equation}
with the following filter:
\begin{equation}
\frac{1}{9}
\begin{bmatrix}
1 & 1 & 1\\
1 & 1 & 1\\
1 & 1 & 1\\
\end{bmatrix}
\end{equation}
is:
\begin{equation}
\frac{1}{9}
\begin{bmatrix}
5 & 4 & 5 & 4 & 5 & 4 \\
4 & 5 & 4 & 5 & 4 & 5 \\
5 & 4 & 5 & 4 & 5 & 4 \\
4 & 5 & 4 & 5 & 4 & 5 \\
5 & 4 & 5 & 4 & 5 & 4 \\
4 & 5 & 4 & 5 & 4 & 5 \\
\end{bmatrix}
\end{equation}
\subsection{} %1.2
This filter works as a blurring filter.
\section{} %2
% make a 3 column table with column names Layer, Output Dimension, Number of Parameters
% CONVk-N(S, P) means a convolutional layer with N kernels of size kxk and stride S and padding P
% padding and stride are 1 if not specified
% POOL-k means a max pooling layer with kernel size kxk
% FC-N means a fully connected layer with N neurons
\begin{table}[h]
\centering
\begin{tabular}{|c|c|c|}
\hline
Layer & Output Dimension & Number of Parameters\\
\hline
Input & $32 \times 32 \times 3$ & 0\\
\hline
CONV3-10 & $32 \times 32 \times 10$ & $3 \times 3 \times 10 + 10 = 280$\\
\hline
ReLU & $32 \times 32 \times 10$ & 0\\
\hline
POOL-2 & $16 \times 16 \times 10$ & 0\\
\hline
CONV3-20(3, 2) & $6 \times 6 \times 20$ & $3 \times 3 \times 20 \times 10 + 20 = 3620$\\
\hline
ReLU & $6 \times 6 \times 20$ & 0\\
\hline
POOL-2 & $3 \times 3 \times 20$ & 0\\
\hline
FLATTEN & 180 & 0\\
\hline
FC-10 & 10 & $180 \times 10 + 10 = 1810$\\
\hline
\end{tabular}
\end{table}
\section{} %3
\subsection{} %3.1
k, b are the parameters of the model.
\subsection{} %3.2
\begin{equation}
\begin{aligned}
	\frac{\partial L}{\partial \omega_1} &= \frac{\partial L}{\partial \hat{y}} \frac{\partial \hat{y}}{\partial \omega_1} \\
	&= (\hat{y} - y) \frac{\partial \hat{y}}{\partial \omega_1} \\
	&= (\hat{y} - y) \nu_1 
\end{aligned}
\end{equation}
\begin{equation}
\begin{aligned}
	\frac{\partial L}{\partial \omega_2} &= \frac{\partial L}{\partial \hat{y}} \frac{\partial \hat{y}}{\partial \omega_2} \\
	&= (\hat{y} - y) \frac{\partial \hat{y}}{\partial \omega_2} \\
	&= (\hat{y} - y) \nu_2
\end{aligned}
\end{equation}
\begin{equation}
	\begin{aligned}
		\frac{\partial L}{\partial a} &= \frac{\partial L}{\partial \hat{y}} \frac{\partial \hat{y}}{\partial a} \\
		&= (\hat{y} - y)
	\end{aligned}
\end{equation}
\subsection{} %3.3
\begin{equation}
	\begin{aligned}
	\frac{\partial L}{\partial z_1} &= \frac{\partial L}{\partial \nu_1} \frac{\partial \nu_1}{\partial z_1} \\
	&= \delta_1 u(z_1 - z_2) u(z_1) = \alpha_1
	\end{aligned}
\end{equation}
\begin{equation}
	\begin{aligned}
	\frac{\partial L}{\partial z_2} &= \frac{\partial L}{\partial \nu_1} \frac{\partial \nu_1}{\partial z_2} + \frac{\partial L}{\partial \nu_2} \frac{\partial \nu_2}{\partial z_2} \\
	&= \delta_1 u(z_2 - z_1) u(z_2) + \delta_2 u(z_2 - z_3) u(z_2) = \alpha_2
	\end{aligned}
\end{equation}
\begin{equation}
	\begin{aligned}
	\frac{\partial L}{\partial z_3} &= \frac{\partial L}{\partial \nu_2} \frac{\partial \nu_2}{\partial z_3} \\
	&= \delta_2 u(z_3 - z_2) u(z_3) = \alpha_3
	\end{aligned}
\end{equation}
Where $ u(x) = \begin{cases} 1 & x > 0 \\ 0 & x \leq 0 \end{cases} $.
\subsection{} %3.4
\begin{equation}
	\begin{aligned}
	\frac{\partial L}{\partial k_1} &= \frac{\partial L}{\partial z_1} \frac{\partial z_1}{\partial k_1} + \frac{\partial L}{\partial z_2} \frac{\partial z_2}{\partial k_1} + \frac{\partial L}{\partial z_3} \frac{\partial z_3}{\partial k_1} \\
	&= \alpha_1 \frac{\partial z_1}{\partial k_1} + \alpha_2 \frac{\partial z_2}{\partial k_1} + \alpha_3 \frac{\partial z_3}{\partial k_1} \\
	&= \alpha_1 x_1 + \alpha_2 x_2 + \alpha_3 x_3 
	\end{aligned}
\end{equation}
\begin{equation}
	\begin{aligned}
	\frac{\partial L}{\partial k_2} &= \frac{\partial L}{\partial z_1} \frac{\partial z_1}{\partial k_2} + \frac{\partial L}{\partial z_2} \frac{\partial z_2}{\partial k_2} + \frac{\partial L}{\partial z_3} \frac{\partial z_3}{\partial k_2} \\
	&= \alpha_1 \frac{\partial z_1}{\partial k_2} + \alpha_2 \frac{\partial z_2}{\partial k_2} + \alpha_3 \frac{\partial z_3}{\partial k_2} \\
	&= \alpha_1 x_2 + \alpha_2 x_3 + \alpha_3 x_4 
	\end{aligned}
\end{equation}
\begin{equation}
	\begin{aligned}
	\frac{\partial L}{\partial k_3} &= \frac{\partial L}{\partial z_1} \frac{\partial z_1}{\partial k_3} + \frac{\partial L}{\partial z_2} \frac{\partial z_2}{\partial k_3} + \frac{\partial L}{\partial z_3} \frac{\partial z_3}{\partial k_3} \\
	&= \alpha_1 \frac{\partial z_1}{\partial k_3} + \alpha_2 \frac{\partial z_2}{\partial k_3} + \alpha_3 \frac{\partial z_3}{\partial k_3} \\
	&= \alpha_1 x_3 + \alpha_2 x_4 + \alpha_3 x_5 
	\end{aligned}
\end{equation}
\begin{equation}
	\begin{aligned}
	\frac{\partial L}{\partial b} &= \frac{\partial L}{\partial z_1} \frac{\partial z_1}{\partial b} + \frac{\partial L}{\partial z_2} \frac{\partial z_2}{\partial b} + \frac{\partial L}{\partial z_3} \frac{\partial z_3}{\partial b} \\
	&= \alpha_1 \frac{\partial z_1}{\partial b} + \alpha_2 \frac{\partial z_2}{\partial b} + \alpha_3 \frac{\partial z_3}{\partial b} \\
	&= \alpha_1 + \alpha_2 + \alpha_3
	\end{aligned}
\end{equation}
\subsection{} %3.5
\begin{equation}
	\begin{aligned}
	\frac{\partial L}{\partial k_j} &= \sum_{i=1}^{m} \frac{\partial L}{\partial z_i} \frac{\partial z_i}{\partial k_j} \\
	&= \sum_{i=1}^{m} \alpha_i \frac{\partial z_i}{\partial k_j} 
	\end{aligned}
\end{equation}
\begin{equation}
	\begin{aligned}
	\frac{\partial z_i}{\partial k_j} &= x_{i+j-1}
	\end{aligned}
\end{equation}
\begin{equation}
	\begin{aligned}
	\frac{\partial L}{\partial k_j} &= \sum_{i=1}^{m} \alpha_i x_{i+j-1} 
	\end{aligned}
\end{equation}
\begin{equation}
	\begin{aligned}
	\frac{\partial L}{\partial b} &= \sum_{i=1}^{m} \alpha_i
	\end{aligned}
\end{equation}
\section{} %4
\subsection{} %4.1
\begin{equation}
\lfloor \frac{205-k}{3} \rfloor + 1 = 66
\end{equation}
$ k = 8, 9, 10$ are all possible values of $ k $.
\subsection{} %4.2
Assuming $ k = 9 $, we have $ 9 \times 9 \times 10 \times 96 + 96 = 8, 064 $ learnable parameters.
\subsection{} %4.3
We have $ 9 \times 9 \times 10 = 810 $ multiplications for each convolution. 
We have $ 65 \times 65 $ of those convolutions. So we have $ 810 \times 65 \times 65 = 4, 985, 250 $ multiplications in total.




\end{document}
